\section{Discussion}
\label{sec:discussion}

\subsection{Confidence as a validation surrogate}

City-scale UBEM cannot wait for per-building ground truth.
Our results suggest that exposing multi-method disagreement and datasource lineage lets practitioners prioritise measurement campaigns and interpret stock-level outputs conservatively.
The negative Spearman correlation between confidence and cross-method spread ($\rho = -0.54$) indicates that ordinal confidence levels carry signal even when absolute accuracy is unknown.

\subsection{Engineering vs.\ research objectives}

Objectives O1--O6 describe \emph{what} we build; hypotheses H1--H6 describe \emph{what} we must prove.
The five-phase methodology (Table~\ref{tab:phases}) links them through validation gates G1--G5.
Pillar P1 (AI accuracy) is strongest today: fine-tuned spatial-SQL models substantially outperform frontier LLMs on CIM PostGIS tasks.
Pillar P2 (semantic fidelity) shows promising VKG coverage with seed OBDA mappings.
Pillar P3 (trustworthiness) depends on completing O5 agents and calibrating confidence weights on additional cities.

\subsection{Limitations}

\textbf{Calibration.}
Confidence weights and tolerances $\tau$ require expert review and empirical tuning per city.

\textbf{Incomplete automation.}
O5 multi-agent VKG generation is not fully implemented; Table~\ref{tab:t4} reports protocol targets.

\textbf{Datasource linkage.}
Not all calculators yet declare \texttt{datasource\_ids}; provenance link rates below 100\% reflect this gap.

\textbf{Case study scope.}
A single pilot city limits generalisation; federated STAC catalogues across municipalities are future work.

\subsection{Future work}

Measured validation where smart-meter or EPC data exist; Brick/SERAF extensions for system-level provenance; learned confidence models replacing hand-tuned rules; and RODI-style OBDA benchmarks for O5.
